%%% Simboli / Symbols %%%

% Inserisci qui tutti i simboli matematici o fisici che utilizzerai.
% Verranno inseriti automaticamente nell'Elenco dei Simboli.
%
% COME USARLI NEL TESTO:
% Usa il comando \gls{etichetta} (es. \gls{sigma} o \gls{alpha}).
% Usare questo metodo è utilissimo: se a metà stesura decidi di cambiare 
% la lettera di una variabile, ti basterà cambiarla qui e si aggiornerà 
% in tutti i capitoli della tesi automaticamente!

% --- Esempi per Ingegneria (Cancellali e inserisci i tuoi) ---

\newglossaryentry{alpha}{
	name={\ensuremath{\alpha}},
	description={Coefficiente di dilatazione termica lineare, in \ensuremath{^\circ\text{C}^{-1}}},
	sort=alpha,
	type=symbols
}

\newglossaryentry{sigma}{
	name={\ensuremath{\sigma}},
	description={Tensione normale applicata a una sezione strutturale, in MPa},
	sort=sigma,
	type=symbols
}

\newglossaryentry{epsilon}{
	name={\ensuremath{\varepsilon}},
	description={Deformazione specifica o assiale (adimensionale)},
	sort=epsilon,
	type=symbols
}

\newglossaryentry{E}{
	name={\ensuremath{E}},
	description={Modulo di elasticità longitudinale (Modulo di Young), in GPa},
	sort=E,
	type=symbols
}